\documentclass[11pt,a4paper]{article}

% ── Codificación, fuentes y lenguaje ───────────────────────────────────────────
\usepackage[utf8]{inputenc}
\usepackage[T1]{fontenc}
\usepackage[default,scale=0.95]{sourcesanspro}
\renewcommand{\familydefault}{\sfdefault}
\usepackage[spanish,es-noshorthands,es-tabla]{babel}

% ── Geometría y espaciado ─────────────────────────────────────────────────────
\usepackage[top=2.2cm, bottom=2.5cm, left=2.2cm, right=2.2cm, headheight=15pt]{geometry}
\usepackage{setspace}
\setstretch{1.18}
\usepackage{parskip}

% ── Matemáticas y unidades ────────────────────────────────────────────────────
\usepackage{amsmath}
\usepackage{amssymb}
\usepackage{siunitx}

% ── Circuitos ─────────────────────────────────────────────────────────────────
\usepackage[european, straightvoltages]{circuitikz}

% ── Tablas ────────────────────────────────────────────────────────────────────
\usepackage{booktabs}
\usepackage{tabularx}
\usepackage{array}
\usepackage{longtable}
\usepackage{colortbl}
\usepackage{enumitem}

% ── Colores, cajas e iconos ───────────────────────────────────────────────────
\usepackage[dvipsnames]{xcolor}
\usepackage{tcolorbox}
\tcbuselibrary{skins, breakable}
\usepackage{fontawesome5}
\usepackage{graphicx}

% ── Listados de código (dark-mode infográfico) ────────────────────────────────
%   Estilo base: fondo oscuro con números de línea. Los bloques lstlisting se
%   envuelven automáticamente en una caja tcolorbox redondeada.
\usepackage{listings}

% ── Secciones con barra de color (infografía) ────────────────────────────────
\usepackage{titlesec}
\titleformat{\section}
  {\normalfont\LARGE\bfseries\color{azulNoche}}
  {\colorbox{cyanNeon}{\color{fondoOscuro}\thesection}}{0.7em}{}
  [\vspace{2pt}\color{cyanNeon}\rule{\linewidth}{1.8pt}]
\titlespacing*{\section}{0pt}{24pt}{10pt}
\titleformat{\subsection}
  {\normalfont\large\bfseries\color{azulNoche}}
  {\textcolor{cyanNeon}{\thesubsection}}{0.7em}{}
  [\vspace{1pt}\color{grisLinea}\rule{\linewidth}{0.6pt}]
\titlespacing*{\subsection}{0pt}{14pt}{6pt}
\titleformat{\subsubsection}
  {\normalfont\normalsize\bfseries\color{azulNoche}}
  {\textcolor{cyanNeon}{\thesubsubsection}}{0.7em}{}
\titlespacing*{\subsubsection}{0pt}{10pt}{4pt}

% ── Cabeceras y pies de página ────────────────────────────────────────────────
\usepackage{fancyhdr}
\usepackage{lastpage}
\pagestyle{fancy}
\fancyhf{}
\renewcommand{\headrulewidth}{0pt}
\renewcommand{\footrulewidth}{0pt}
\fancyfoot[C]{%
  \begin{minipage}{\linewidth}
    \vspace{2pt}
    \color{cyanNeon}\rule{\linewidth}{1.2pt}\\[2pt]
    \centering\color{grisTexto}\footnotesize
    Página \thepage\ de \pageref{LastPage}
  \end{minipage}%
}

% ── Hiperenlaces ──────────────────────────────────────────────────────────────
\usepackage[colorlinks=true, linkcolor=azulNoche, urlcolor=cyanNeon]{hyperref}
\sloppy
\emergencystretch 3em

% ── Paleta de colores — tecnológico dark-mode ──────────────────────────────────
\definecolor{fondoOscuro}{HTML}{0D1B2A}     % Dark navy — fondo banda título
\definecolor{verdeTurquesa}{HTML}{004D40}    % Verde turquesa — bandas de marca
\definecolor{azulNoche}{HTML}{1B3A6B}       % Midnight blue — secciones, marcos
\definecolor{azulMedio}{HTML}{2A5298}       % Azul medio — variante de acento
\definecolor{cyanNeon}{HTML}{00C8E0}        % Cyan eléctrico — acentos primarios
\definecolor{verdeSignal}{HTML}{00B887}     % Verde señal — fortalezas, éxito
\definecolor{naranjaVivo}{HTML}{FF7043}     % Naranja vivo — mejoras, advertencias
\definecolor{rojoAlerta}{HTML}{E53935}      % Rojo alerta — errores
\definecolor{grisPapel}{HTML}{F4F6FA}       % Fondo tarjetas claras
\definecolor{grisLinea}{HTML}{C8D0E0}       % Bordes sutiles
\definecolor{grisTexto}{HTML}{6B7A99}       % Texto secundario
\definecolor{amarilloNota}{HTML}{FFB300}    % Notas / advertencias doradas

% Colores específicos para bloques de código (dark-mode)
\definecolor{codigoFondo}{HTML}{1E2D3D}      % Fondo del bloque de código
\definecolor{codigoComentario}{HTML}{7A8FA6} % Comentarios
\definecolor{codigoCadena}{HTML}{FF9D5C}     % Cadenas / strings
\definecolor{codigoKeyword}{HTML}{00C8E0}    % Palabras clave
\definecolor{codigoNumero}{HTML}{00E5A0}     % Números

% ── Banda de título: portada infográfica ──────────────────────────────────────
%   Diseño en 2 capas: banda de color (por defecto fondo oscuro) + franja de
%   acento cyan abajo. Uso: \bandaTitulo[color]{título}{subtítulo}.
%   El icono se puede cambiar con \renewcommand{\iconoBanda}{...} (FontAwesome).
\newcommand{\iconoBanda}{microchip}
\newcommand{\bandaTitulo}[3][fondoOscuro]{%
  \begin{tcolorbox}[
    enhanced,
    colback=#1, colframe=#1,
    boxrule=0pt, arc=8pt, outer arc=8pt,
    width=\linewidth,
    top=18pt, bottom=6pt, left=14pt, right=14pt,
    borderline south={3.5pt}{0pt}{cyanNeon},
  ]
    \begin{minipage}[c]{0.07\linewidth}
      \centering
      \textcolor{cyanNeon}{\fontsize{30}{30}\selectfont\faIcon{\iconoBanda}}
    \end{minipage}%
    \hspace{8pt}%
    \begin{minipage}[c]{0.88\linewidth}
      {\fontsize{20}{24}\selectfont\bfseries\color{white}#2}\par\vspace{5pt}%
      \textcolor{cyanNeon!80!white}{\small\faIcon{angle-right}~#3}%
    \end{minipage}
    \vspace{4pt}
  \end{tcolorbox}%
  \vspace{8pt}%
}

% ── Tarjetas de datos (infografía) ────────────────────────────────────────────
\tcbset{
  tarjetaDato/.style={
    enhanced, breakable,
    arc=6pt, outer arc=6pt,
    colback=grisPapel, colframe=azulNoche,
    boxrule=1pt,
    fonttitle=\bfseries\small, coltitle=white,
    attach boxed title to top left={yshift=-3mm, xshift=6mm},
    boxed title style={
      colback=azulNoche, colframe=azulNoche,
      arc=4pt, boxrule=0pt,
    },
    drop fuzzy shadow=azulNoche!30!white,
    top=8pt, bottom=8pt, left=10pt, right=10pt,
  },
  tarjetaIcono/.style={
    enhanced, breakable,
    arc=6pt, outer arc=6pt,
    colback=white, colframe=grisLinea, boxrule=0.8pt,
    fonttitle=\bfseries\small, coltitle=azulNoche,
    borderline west={3pt}{0pt}{cyanNeon},
    drop fuzzy shadow=azulNoche!25!white,
    top=6pt, bottom=6pt, left=10pt, right=10pt,
  },
  cajaTitulo/.style={
    enhanced, breakable,
    arc=6pt, outer arc=6pt,
    colback=azulNoche!10!grisPapel, colframe=azulNoche,
    boxrule=1pt,
    fonttitle=\bfseries, coltitle=white,
    attach boxed title to top left={yshift=-3mm, xshift=6mm},
    boxed title style={
      colback=azulNoche, colframe=azulNoche,
      arc=4pt, boxrule=0pt,
    },
    drop fuzzy shadow=azulNoche!25!white,
    top=8pt, bottom=8pt, left=10pt, right=10pt,
  },
  cajaContenido/.style={
    enhanced, breakable,
    arc=5pt, outer arc=5pt,
    colback=grisPapel, colframe=grisLinea,
    boxrule=0.8pt,
    fonttitle=\bfseries\small, coltitle=azulNoche,
    top=6pt, bottom=6pt, left=10pt, right=10pt,
  },
  cajaFortaleza/.style={
    enhanced, breakable,
    arc=6pt, outer arc=6pt,
    colback=verdeSignal!8!white, colframe=verdeSignal,
    boxrule=1pt,
    borderline west={4pt}{0pt}{verdeSignal},
    fonttitle=\bfseries\small, coltitle=verdeSignal!70!black,
    drop fuzzy shadow=verdeSignal!20!white,
    top=6pt, bottom=6pt, left=10pt, right=10pt,
  },
  cajaMejora/.style={
    enhanced, breakable,
    arc=6pt, outer arc=6pt,
    colback=naranjaVivo!8!white, colframe=naranjaVivo,
    boxrule=1pt,
    borderline west={4pt}{0pt}{naranjaVivo},
    fonttitle=\bfseries\small, coltitle=naranjaVivo!80!black,
    drop fuzzy shadow=naranjaVivo!20!white,
    top=6pt, bottom=6pt, left=10pt, right=10pt,
  },
  cajaRecomendacion/.style={
    enhanced, breakable,
    arc=6pt, outer arc=6pt,
    colback=cyanNeon!6!white, colframe=cyanNeon!60!azulNoche,
    boxrule=1pt,
    borderline west={4pt}{0pt}{cyanNeon},
    fonttitle=\bfseries\small, coltitle=azulNoche,
    drop fuzzy shadow=azulNoche!20!white,
    top=6pt, bottom=6pt, left=10pt, right=10pt,
  },
}

% ── Ícono + texto (encabezados de sección y elementos) ────────────────────────
\newcommand{\iconotexto}[2]{\textcolor{cyanNeon}{\faIcon{#1}}~#2}

% ── Listados de código: estilo dark + auto-envoltura ─────────────────────────
\lstdefinestyle{estiloCodigo}{
    backgroundcolor=\color{codigoFondo},
    basicstyle=\ttfamily\small\color{white},
    commentstyle=\color{codigoComentario}\itshape,
    stringstyle=\color{codigoCadena},
    keywordstyle=\color{codigoKeyword}\bfseries,
    numberstyle=\tiny\color{codigoComentario},
    numbers=left,
    numbersep=10pt,
    showstringspaces=false,
    breaklines=true,
    breakatwhitespace=false,
    tabsize=4,
    frame=none,
    captionpos=b,
    aboveskip=6pt,
    belowskip=4pt,
    xleftmargin=14pt,
    framexleftmargin=14pt,
    literate=
      {á}{{\'a}}1 {é}{{\'e}}1 {í}{{\'i}}1 {ó}{{\'o}}1 {ú}{{\'u}}1
      {Á}{{\'A}}1 {É}{{\'E}}1 {Í}{{\'I}}1 {Ó}{{\'O}}1 {Ú}{{\'U}}1
      {ñ}{{\~n}}1 {Ñ}{{\~N}}1 {ü}{{\"u}}1 {Ü}{{\"U}}1
      {¿}{{?`}}1 {¡}{{!`}}1
}
\lstdefinestyle{estiloPython}{
    style=estiloCodigo,
    language=Python,
}
\lstset{style=estiloCodigo}
\BeforeBeginEnvironment{lstlisting}{%
  \begin{tcolorbox}[
    enhanced, breakable,
    arc=6pt, outer arc=6pt,
    colback=codigoFondo, colframe=azulNoche,
    boxrule=1pt,
    drop fuzzy shadow=azulNoche!25!white,
    top=2pt, bottom=2pt, left=0pt, right=0pt,
  ]%
}
\AfterEndEnvironment{lstlisting}{\end{tcolorbox}}

% ── Cajas de contenido temáticas ──────────────────────────────────────────────
\newtcolorbox{cajaRecuerda}{
    enhanced, breakable,
    arc=6pt, outer arc=6pt,
    colback=azulNoche!10!grisPapel,
    colframe=azulNoche, boxrule=1pt,
    borderline west={4pt}{0pt}{cyanNeon},
    fonttitle=\bfseries\small\color{white},
    title={\faIcon{info-circle}~Recuerda},
    attach boxed title to top left={yshift=-3mm, xshift=6mm},
    boxed title style={colback=azulNoche, arc=4pt, boxrule=0pt},
    drop fuzzy shadow=azulNoche!25!white,
    left=10pt, right=10pt, top=8pt, bottom=8pt,
}

\newtcolorbox{cajaConcepto}{
    enhanced, breakable,
    arc=6pt, outer arc=6pt,
    colback=verdeSignal!8!white,
    colframe=verdeSignal, boxrule=1pt,
    borderline west={4pt}{0pt}{verdeSignal},
    fonttitle=\bfseries\small\color{white},
    title={\faIcon{lightbulb}~Concepto clave},
    attach boxed title to top left={yshift=-3mm, xshift=6mm},
    boxed title style={colback=verdeSignal!80!black, arc=4pt, boxrule=0pt},
    drop fuzzy shadow=verdeSignal!20!white,
    left=10pt, right=10pt, top=8pt, bottom=8pt,
}

\newtcolorbox{cajaEjemplo}{
    enhanced, breakable,
    arc=6pt, outer arc=6pt,
    colback=cyanNeon!5!white,
    colframe=cyanNeon!70!azulNoche, boxrule=1pt,
    borderline west={4pt}{0pt}{cyanNeon},
    fonttitle=\bfseries\small\color{fondoOscuro},
    title={\faIcon{code}~Ejemplo},
    attach boxed title to top left={yshift=-3mm, xshift=6mm},
    boxed title style={colback=cyanNeon, arc=4pt, boxrule=0pt},
    drop fuzzy shadow=azulNoche!20!white,
    left=10pt, right=10pt, top=8pt, bottom=8pt,
}

\newtcolorbox{cajaEjercicio}{
    enhanced, breakable,
    arc=6pt, outer arc=6pt,
    colback=azulMedio!7!white,
    colframe=azulMedio, boxrule=1pt,
    borderline west={4pt}{0pt}{azulMedio},
    fonttitle=\bfseries\small\color{white},
    title={\faIcon{pencil-alt}~Ejercicio},
    attach boxed title to top left={yshift=-3mm, xshift=6mm},
    boxed title style={colback=azulMedio, arc=4pt, boxrule=0pt},
    drop fuzzy shadow=azulNoche!20!white,
    left=10pt, right=10pt, top=8pt, bottom=8pt,
}

\newtcolorbox{cajaReto}{
    enhanced, breakable,
    arc=6pt, outer arc=6pt,
    colback=naranjaVivo!6!white,
    colframe=naranjaVivo, boxrule=1pt,
    borderline west={4pt}{0pt}{naranjaVivo},
    fonttitle=\bfseries\small\color{white},
    title={\faIcon{fire}~Reto},
    attach boxed title to top left={yshift=-3mm, xshift=6mm},
    boxed title style={colback=naranjaVivo, arc=4pt, boxrule=0pt},
    drop fuzzy shadow=naranjaVivo!20!white,
    left=10pt, right=10pt, top=8pt, bottom=8pt,
}

% ── Indicadores de dificultad ─────────────────────────────────────────────────
\newcommand{\facil}{\textcolor{verdeSignal}{\faIcon{circle}\,\textbf{Fácil}}}
\newcommand{\intermedio}{\textcolor{naranjaVivo}{\faIcon{circle}\,\textbf{Intermedio}}}
\newcommand{\dificil}{\textcolor{rojoAlerta}{\faIcon{circle}\,\textbf{Difícil}}}

% ─────────────────────────────────────────────────────────────────────────────


% ── Cabecera específica de este documento ──────────────────────────────────────
\fancyhead[L]{\small\color{azulNoche}\textbf{Workstation para Agentes IA --- AM5}}
\fancyhead[R]{\small\color{grisTexto}\itshape B650 Gaming X AX V2}
\renewcommand{\headrulewidth}{0.6pt}
\renewcommand{\headrule}{\color{cyanNeon}\hrule width\headwidth height 0.6pt}
\renewcommand{\iconoBanda}{desktop}


\begin{document}

% ══ PORTADA INFográfica ══════════════════════════════════════════════════════
\bandaTitulo{Lista de Compras: Workstation para Agentes IA}{Gigabyte B650 Gaming X AX V2 --- Plataforma AM5, 100\,\% vía APIs (sin GPU dedicada)}

\begin{center}
  \vspace{2pt}
  {\large \textbf{Autor:} Prof. CERO \hfill \textbf{Fecha:} \today} \\[6pt]
  \color{cyanNeon}\rule{\linewidth}{0.4pt}
\end{center}

\tableofcontents
\newpage

% ══ 1. RESUMEN ══════════════════════════════════════════════════════════════
\section[Resumen]{\iconotexto{info-circle}{Resumen}}

Esta es la lista de compras completa para ensamblar una \textbf{estación de trabajo dedicada al desarrollo de Agentes de IA}, orientada a un flujo \textbf{100\,\% basado en APIs} (OpenRouter, Gemini, Groq), sin necesidad de GPU dedicada para inferencia local.

La base es la placa \textbf{Gigabyte B650 Gaming X AX V2}, sobre la plataforma \textbf{AM5} de AMD. El presupuesto de la GPU se redistribuye en \textbf{más RAM y almacenamiento}, que son los cuellos de botella reales de un flujo de desarrollo de agentes (ChromaDB, embeddings, múltiples procesos Python, servidores MCP y orquestación en paralelo).

\begin{tcolorbox}[cajaFortaleza, title={\faIcon{thumbs-up}~Concepto de la build}]
El cómputo pesado (inferencia de modelos) lo realiza el LLM remoto. Esta máquina es la \textbf{orquestadora}: corre Python, LangChain, ChromaDB, Docker, LaTeX y múltiples terminales en paralelo sin fricción.
\end{tcolorbox}

% ══ 2. LISTA DE COMPONENTES ══════════════════════════════════════════════════
\section[Lista de Componentes]{\iconotexto{shopping-cart}{Lista de Componentes}}

\begin{longtable}{p{3.2cm} p{6.2cm} p{2.2cm} p{0.6cm}}
\toprule
\textbf{Componente} & \textbf{Modelo recomendado} & \textbf{Precio aprox.} & \textbf{Link} \\
\midrule
\endhead
\bottomrule
\endlastfoot

\textbf{Motherboard} & Gigabyte B650 Gaming X AX V2 (ATX, AM5, WiFi 6E, 2.5G LAN) & \$160 & \ref{link-mb} \\
\textbf{CPU} & AMD Ryzen 7 7700X (8C/16T, Zen 4) & \$280 & \ref{link-cpu} \\
\textbf{RAM} & G.Skill Flare X5 32\,GB (2$\times$16\,GB) DDR5-5600 -- \textbf{x2 para 64\,GB} & \$180 & \ref{link-ram} \\
\textbf{SSD 1 (OS)} & WD Black SN770 1\,TB NVMe PCIe 4.0 & \$80 & \ref{link-ssd1} \\
\textbf{SSD 2 (Datos)} & WD Black SN770 2\,TB NVMe PCIe 4.0 & \$120 & \ref{link-ssd2} \\
\textbf{PSU} & Corsair CX650 650\,W 80\,Plus Bronze & \$60 & \ref{link-psu} \\
\textbf{Cooler} & Thermalright Peerless Assassin 120 SE & \$35 & \ref{link-cooler} \\
\textbf{Case} & Fractal Design Meshify C o similar con frontal mesh & \$80 & \ref{link-case} \\
\bottomrule
\end{longtable}

\begin{tcolorbox}[tarjetaDato, title={\faIcon{coins}~Costo total estimado}]
\centering
{\fontsize{28}{34}\selectfont\bfseries\color{azulNoche}\$995 USD}\par\vspace{4pt}
{\color{grisTexto}\small Aprox., sin GPU dedicada (mercado actual)}
\end{tcolorbox}

% ══ 3. DETALLE POR COMPONENTE CON LINKS ══════════════════════════════════════
\section[Detalle y Compras]{\iconotexto{link}{Detalle de Cada Componente con Enlaces de Compra}}

\subsection{Motherboard --- Gigabyte B650 Gaming X AX V2}
\label{link-mb}
\begin{tcolorbox}[cajaRecomendacion, title={\faIcon{memory}~Placa base}]
\textbf{Elección central.} Plataforma AM5, formato ATX, WiFi 6E, LAN 2.5\,GbE, 2$\times$ M.2 NVMe PCIe 4.0 y VRM de 12$+$2 fases. Suficiente para una estación de trabajo sin tarjeta gráfica dedicada.

Compra: \url{https://www.amazon.com/s?k=Gigabyte+B650+Gaming+X+AX+V2}
Alternativa: \url{https://pcpartpicker.com/search/?q=B650+Gaming+X+AX+V2}
\end{tcolorbox}

\subsection{CPU --- AMD Ryzen 7 7700X}
\label{link-cpu}
\begin{tcolorbox}[cajaRecomendacion, title={\faIcon{microchip}~Procesador}]
\textbf{8 n\'ucleos / 16 hilos}, Zen 4, TDP de 105\,W. Potencia suficiente para orquestar agentes, compilar LaTeX y correr Docker y servidores MCP en paralelo. Si el presupuesto ajusta, el \textbf{Ryzen 7 7700} (65\,W) es una opción m\'as eficiente.

Compra: \url{https://www.amazon.com/s?k=AMD+Ryzen+7+7700X}
Alternativa: \url{https://pcpartpicker.com/search/?q=Ryzen+7+7700X}
\end{tcolorbox}

\subsection{RAM --- 64\,GB DDR5 (2 paquetes de 32\,GB)}
\label{link-ram}
\begin{tcolorbox}[cajaRecomendacion, title={\faIcon{microchip}~Memoria}]
\textbf{64\,GB (4$\times$16\,GB o 2$\times$32\,GB)} DDR5-5600 optimizada para AMD (EXPO). Se recomienda un kit de \textbf{2$\times$32\,GB} para dejar libertad de ampliación a 128\,GB. Crítica para ChromaDB, embeddings y m\'ultiples procesos.

Compra: \url{https://www.amazon.com/s?k=DDR5+5600+AMD+EXPO+32GB+kit}
Alternativa: \url{https://pcpartpicker.com/products/memory/}
\end{tcolorbox}

\subsection{Almacenamiento --- Doble NVMe}
\label{link-ssd1}\label{link-ssd2}
\begin{tcolorbox}[cajaRecomendacion, title={\faIcon{hdd}~SSD}]
\textbf{SSD 1:} 1\,TB NVMe PCIe 4.0 para el sistema operativo, aplicaciones y entornos conda. \textbf{SSD 2:} 2\,TB NVMe PCIe 4.0 para proyectos, datasets y la base `chroma\_db`. Separarlos evita colisiones de E/S.

Compra (1\,TB): \url{https://www.amazon.com/s?k=WD+Black+SN770+1TB}
Compra (2\,TB): \url{https://www.amazon.com/s?k=WD+Black+SN770+2TB}
\end{tcolorbox}

\subsection{Fuente de Poder --- 650\,W 80\,Plus}
\label{link-psu}
\begin{tcolorbox}[cajaRecomendacion, title={\faIcon{bolt}~PSU}]
\textbf{650\,W 80\,Plus Bronze} es m\'as que suficiente sin GPU dedicada. Deja margen si m\'as adelante se a\~nade una aceleradora de gama media (hasta 250\,W).

Compra: \url{https://www.amazon.com/s?k=Corsair+CX650}
Alternativa: \url{https://pcpartpicker.com/products/power-supply/}
\end{tcolorbox}

\subsection{Disipador --- Torre de aire eficiente}
\label{link-cooler}
\begin{tcolorbox}[cajaRecomendacion, title={\faIcon{fan}~Cooler}]
Una torre de aire de doble ventilador (tipo \textbf{Thermalright Peerless Assassin 120 SE}) rinde mejor que el disipador de stock para cargas sostenidas de compilaci\'on, manteniendo el Ryzen 7 fresco y silencioso.

Compra: \url{https://www.amazon.com/s?k=Thermalright+Peerless+Assassin+120+SE}
\end{tcolorbox}

\subsection{Gabinete --- Frontal Mesh}
\label{link-case}
\begin{tcolorbox}[cajaRecomendacion, title={\faIcon{box-open}~Case}]
\textbf{Formato ATX} con frontal de malla perforada y buen flujo de aire. No se necesita un gabinete ex\'otico; un chasis mesh medio basta para disipar la baja carga t\'ermica de esta build.

Compra: \url{https://www.amazon.com/s?k=ATX+mesh+case}
Alternativa: \url{https://pcpartpicker.com/products/case/}
\end{tcolorbox}

% ══ 4. LO QUE NO HACE FALTA COMPRAR ══════════════════════════════════════════
\section[Qu\'e NO comprar]{\iconotexto{ban}{Qu\'e NO hace falta comprar}}

\begin{tcolorbox}[cajaMejora, title={\faIcon{exclamation-circle}~Ahorros conscientes}]
\begin{itemize}
    \item \textbf{GPU dedicada:} No es necesaria para un flujo 100\,\% API. Ahorro de \$300--\$800.
    \item \textbf{PSU de alta potencia (850\,W+):} Sobra. Ahorro de \$40+.
    \item \textbf{Enfriamiento l\'iquido:} Innecesario a 105\,W. Ahorro de \$80+.
    \item \textbf{PCIe 5.0 en GPU:} B650 no lo ofrece y no se necesita para orquestaci\'on.
\end{itemize}
\end{tcolorbox}

% ══ 5. VEREDICTO ════════════════════════════════════════════════════════════
\section[Veredicto]{\iconotexto{flag-checkered}{Veredicto Final}}

\begin{tcolorbox}[cajaFortaleza, title={\faIcon{thumbs-up}~Conclusi\'on}]
Esta build entrega la mejor relaci\'on rendimiento-precio para el desarrollo de Agentes de IA v\'ia API: CPU de 8 n\'ucleos, 64\,GB de RAM DDR5 y doble NVMe, todo sobre la s\'olida plataforma AM5 de la \textbf{B650 Gaming X AX V2}, sin gastar un d\'olar de m\'as en hardware que no se aprovechar\'ia.

El \'unico upgrade futuro genuino ser\'ia a\~nadir una GPU solo si posteriormente decides correr modelos peque\~nos localmente.
\end{tcolorbox}

\end{document}
